\documentclass[11pt]{article}
\usepackage[margin=1in]{geometry}
\usepackage{ebgaramond}
\usepackage{setspace}
\setstretch{1.15}
\usepackage{graphicx}
\usepackage{grffile}
\usepackage{booktabs}
\usepackage{amsmath}
\usepackage{amssymb}
\usepackage{caption}
\usepackage{ragged2e}
\usepackage{xcolor}
\usepackage[colorlinks=true, linkcolor=blue!70!black, citecolor=blue!70!black, urlcolor=blue!70!black]{hyperref}
\urlstyle{same}
\renewcommand{\ttdefault}{cmtt}

% Caption styling: figure title above figure, unbolded
\captionsetup[figure]{labelfont=normalfont,textfont=normalfont,position=top,skip=6pt}

% Figure note styling: reduced font size (2 sizes smaller, italic)
\newcommand{\figurenotesize}{\scriptsize}

% Section numbering format:
% For major section titles (\section), outdent the number so the title text aligns with the start of the body text
\makeatletter
\def\@seccntformat#1{%
  \ifcsname format@#1\endcsname
    \csname format@#1\endcsname
  \else
    \csname the#1\endcsname.\quad
  \fi
}
\def\format@section{\llap{\thesection.\hspace{0.8em}}}
\makeatother

% Graphic search paths into R - Results/R1.1 - Figures/
\graphicspath{
  {../R - Results/R1.1 - Figures/R1.1.2 - T - Three Step Text Analysis Pipeline/R1.1.2.1 - T1 - Text Summarization/}
  {../../R - Results/R1.1 - Figures/R1.1.2 - T - Three Step Text Analysis Pipeline/R1.1.2.1 - T1 - Text Summarization/}
  {../R - Results/R1.1 - Figures/R1.1.2 - T - Three Step Text Analysis Pipeline/R1.1.2.2 - T2 - Topic Extraction and Stance Analysis/}
  {../../R - Results/R1.1 - Figures/R1.1.2 - T - Three Step Text Analysis Pipeline/R1.1.2.2 - T2 - Topic Extraction and Stance Analysis/}
  {../R - Results/R1.1 - Figures/R1.1.2 - T - Three Step Text Analysis Pipeline/R1.1.2.3 - T3 - Topic Relevance and Corporate Disclosure/}
  {../../R - Results/R1.1 - Figures/R1.1.2 - T - Three Step Text Analysis Pipeline/R1.1.2.3 - T3 - Topic Relevance and Corporate Disclosure/}
  {../R - Results/R1.1 - Figures/R1.1.3 - X1 - Example of Research Enabled by LLM Pipeline/R1.1.3.1 - Figure 4 - Firm Topic Relevance and Privacy Patent Changes Pre-Post 2016/}
  {../../R - Results/R1.1 - Figures/R1.1.3 - X1 - Example of Research Enabled by LLM Pipeline/R1.1.3.1 - Figure 4 - Firm Topic Relevance and Privacy Patent Changes Pre-Post 2016/}
  {../R - Results/R1.1 - Figures/R1.1.3 - X1 - Example of Research Enabled by LLM Pipeline/R1.1.3.2 - Figure 5 - Long-Term Trends and Firm Concentration in Privacy Patenting/}
  {../../R - Results/R1.1 - Figures/R1.1.3 - X1 - Example of Research Enabled by LLM Pipeline/R1.1.3.2 - Figure 5 - Long-Term Trends and Firm Concentration in Privacy Patenting/}
  {../R - Results/R1.1 - Figures/R1.1.3 - X1 - Example of Research Enabled by LLM Pipeline/R1.1.3.3 - Figure 6 - Annual Event-Study Estimates by Topic Relevance Quartile/}
  {../../R - Results/R1.1 - Figures/R1.1.3 - X1 - Example of Research Enabled by LLM Pipeline/R1.1.3.3 - Figure 6 - Annual Event-Study Estimates by Topic Relevance Quartile/}
}

% Adjust list spacing
\def\labelitemi{--}

\begin{document}

\begin{center}
    {\LARGE \textsc{Opening the Black Box of Large Language Models in Text-based Measurement: Staged Pipelines with Verifiable Checkpoints}\par}
    \vspace{1.2em}
    {\normalsize\bfseries Milan Miric, Albert Roh, Nan Jia\par}
    \vspace{0.2em}
    {\small\itshape University of Southern California, Marshall School of Business\par}
    \vspace{0.8em}
    {\footnotesize\textbf{Proposal for Special Issue:} AI-Enabled Frontiers in Organizational Science (\textit{Organization Science})\par}
    \vspace{0.3em}
    {\footnotesize\textbf{Repository:} \url{https://github.com/HC-Geopolitics-Innovation/LLM-Text-Analysis-and-Synthesis}\par}
\end{center}

\vspace{1.5em}

\section{The Data Opportunity and Black-box Risk in LLM-based Measurement: A Staged Pipeline Approach}

Researchers increasingly use large language models (LLMs) to construct variables from text (Gilardi et al. 2023, Rathje et al. 2024, Asirvatham et al. 2026, Dell and Rambachan 2026). The appeal is clear: long noisy documents which have typically been unsuitable for processing with traditional techniques such as topic models or embeddings (Grimmer and Stewart 2013, Hannigan et al. 2019, Dell 2025), such as 10-K filings and policy reports, can be processed by LLMs (Kim et al. 2023, Jha et al. 2024, de Kok 2025, Benoit et al. 2026). This opens research contexts that were previously impractical to study at very large scale. However, when LLMs are used as an end to end solution: by providing a model with entire documents and asking it to return a numerical quantity, the LLM becomes a black box in that the final output is the only observable checkpoint. Researchers cannot see which parts of a document a measure rests on, what was omitted, or how judgments are made at intermediate stages. In addition, the outcome is sensitive to unobserved choices, making it difficult to reproduce.

This opacity creates three risks. (1) \textit{Validation.} Validating only the final output shows how often a measure is accurate in specific cases, but provides no systematic indication why it may have not been accurate. These aggregate measures can mask systematic errors: Carlson and Burbano (2026) show how these measures need to be validated by creating a heterogeneous set of approaches and then validating the outcomes on those approaches. (2) \textit{Reproducibility.} Rerunning an end-to-end process may not return the same measure when proprietary models are updated or retired; and because the process produces no intermediate outputs, other researchers cannot leverage the intermediate outputs generated by that process, and instead must re-create the entire effort. (3) \textit{Heterogeneity within documents.} An end-to-end measure collapses each document into a single output, concealing variation inside it. Researchers do not have control over which portion of a text document is driving the final outputs of the large language models, or whether for some parts of the model, the model is creating summaries and not processing these document equally. Because these measures enter subsequent analyses as variables, the risk is not merely noisy measurement but biased analysis: when errors are systematic, a finding may reflect how a measure was produced rather than the phenomenon it is meant to capture. One might expect these risks to fade as models gain longer context windows and stronger reasoning. However, many of these concerns will remain. A more capable end-to-end model will only expose the final output, but the concerns will emerge from opacity and a lack of control over design choices, rather than from a limited capability.

Recent work has acknowledged and begun to devise methodologies to overcome these risks. Carlson and Burbano (2026) propose a way of mitigating these concerns by systematic sensitivity analysis, including generating families of prompt variations, running them across models, and reporting the range of resulting downstream estimates. This approach reveals how much conclusions depend on implementation choices, but also focuses on a single albeit common task of textual classification. Processing long text documents and extracting specific insights nests a text classification choice within a black box of other data parsing, processing and synthesis functions. Using large language models for these other more complex textual analysis tasks, creates more opportunities for the models to make choices. This is part of a broader literature which has emerged showing the idiosyncrasies of how these models can make individual choices, but when combined as part of a carefully crafted workflow can mimic at scale or improve what human workers would have achieved (Csaszar et al. 2024, Doshi et al. 2025, Dell'Acqua et al. 2026). These studies have demonstrated important AI use in specific tasks, many of which could be easily integrated into an AI pipeline. However, the LLM processing of large information rich text documents which could enable researchers to answer new and important questions, requires finding a way to integrate and apply these methodologies as part of a pipeline which could in turn avoid introducing systematic noise and bias through black box algorithms. We propose a complementary approach.

Instead of using an end-to-end approach (e.g. using an LLM to parse, summarize and process an entire document) and varying the input (prompt text), we propose opening up the black-box algorithm and decomposing LLM-based measurement generation into a pipeline of discrete stages, each producing an inspectable intermediate output. This design addresses each risk. For validation, each stage can be checked against its own benchmark, allowing researchers to detect systematic error and locate where it arises. For reproducibility, a change in model, or a new construct definition, is confined to one stage, so researchers can reuse the outputs of earlier stages and revalidate only what changed; the same structure lets sensitivity analysis vary one stage at a time and attribute any change in results to it. For heterogeneity, intermediate outputs are retained for individual passages rather than whole documents, so aggregate measures can be traced back to their components and re-aggregated by contributor, section, or other conditions of interest.

Decomposing the process also allows each stage to use the tool best suited to it. LLMs excel at two types of tasks: (1) Synthesizing text, that is, reading long, noisy passages and condensing them into short, information-dense outputs such as summaries or statements of position; and (2) Generating code, that is, producing machine-readable instructions that can be systematically executed and tested. We use LLMs for synthesis and rely on established methods, such as embeddings and similarity measures, where they have proven routines, validation standards, and verifiable properties, applying them to condensed or sentence-level text rather than to whole documents. The pipeline thus combines the core strengths of LLMs with the known properties of established tools rather than replacing one with the other. Code generation also makes the pipeline easier to extend to other use cases: because the pipeline's code is available in a public repository, other researchers can use coding agents to adapt it to their own data and questions (Appendix D).

We demonstrate the application of this approach in a setting relevant for organization researchers: how congressional debate over regulating a technology relates to changes in the innovation trajectories of firms exposed to it. Studying this question requires understanding and quantifying the nature of the policy debate which exists and how it might be leveraged by researchers. This also requires measuring the extent to which the issues being discussed in the debate represent threats or opportunities for firms.

We focus on the policy debate over user data privacy that emerged in the mid 2010's related to technology companies collecting customer data, the Cambridge Analytica scandal related to the use of Facebook data by third party companies, and document the relationship between the debate of policy makers which coincided with shifts in the innovation policy of firms. Studying this question requires identifying which firms were likely impacted by this debate, so they can be compared with firms that credibly were not. That information is embedded in two kinds of long, noisy documents: congressional hearing transcripts, which average roughly 40,000 words and alternate between testimony and questioning from many participants, and firms' 10-K filings. The setting puts all three risks at stake: Because the exposure measure assigns firms to treatment and control groups, systematic error in it would bias the estimated effect; because other researchers may want to extend the measure to studying other policies, it must be reproducible in parts; and because a committee's stance may reflect its members broadly or only a few vocal participants, heterogeneity within hearings matters for interpretation.

\section{Our Framework: A Staged Pipeline for LLM-based Measurement}

We develop a framework that illustrates how LLMs can be used within a data pipeline without ceding full control to them. After data preparation, the pipeline proceeds through three stages (text summarization, stance extraction, and relevance to firms' business areas), each composed of discrete steps that produce intermediate outputs. Each output can be validated, reused, and disaggregated, addressing the three risks above. Researchers can adapt these stages to extract different insights or apply them to data of their choice.

\subsection{Overview of Data Pipeline and Processing Stages}

\begin{itemize}
    \item \textbf{Stage 1 -- Data Preparation:} Assembling and structuring the corpora before any LLM is applied. We parse congressional hearing transcripts from the U.S. Government Publishing Office into speaker-attributed turns, extract business descriptions (Item 1) from S\&P 500 firms' 10-K filings, and retrieve independent C-SPAN hearing abstracts as a validation benchmark. Defects such as undigitized transcripts are flagged rather than dropped, so every exclusion can be inspected.
    \item \textbf{Stage 2 -- Text Summarization:} Taking large, complex textual documents that contain considerable noise alongside embedded data, partitioning them into chunks, and using LLMs to generate, synthesize, and aggregate textual summaries. We demonstrate this approach on congressional hearing transcripts (see Appendix B for data sources and processing), summarizing long and convoluted documents into concise summaries that can subsequently be used for text classification, topic modeling, or semantic distance estimation in more reliable ways.
    \item \textbf{Stage 3 -- Stance Extraction:} Using LLMs to identify the positions of different actors (e.g., members of Congress, firms, government agencies) on specific topics, allowing researchers to distill the nuances of arguments and positions that each individual takes on an issue. We apply this to congressional hearings to identify the stances of members of Congress on key issues, classified as support, concern, or neutral, in a format that can be utilized in subsequent research.
    \item \textbf{Stage 4 -- Debate Relevance to Firm Business Areas:} Using the distilled stances of individuals within the debate, we create a synthesis of the policy debate. This synthesis, or an embedding of the underlying statements, can then be compared to the business descriptions of firms to identify which firms are exposed to the debate. In our empirical context, this means framing the key issues that policymakers are debating and attempting to regulate. We then compare this to the business descriptions (Item 1) in firms' 10-K filings. This allows us to identify and study firms that are exposed to the debate versus those that are not.
\end{itemize}

\subsection{Adapting the Pipeline for Other Research Questions}

Because each stage is modular and its code is available in our repository, researchers can reuse the pipeline in part or in full. For example, they can apply the same stages to hearings of other committees or to other policy topics, or replace the stance-extraction stage to measure a different construct while keeping data preparation and summarization fixed.

Furthermore, LLMs' ability to generate code lowers the cost of such adaptation: these models also excel at writing code that is automatically tested, verified, and validated. By providing our pipeline's code and methodologies in a GitHub repository where the codebase is available, researchers can use their own coding agents to either improve upon our methodology or apply this methodology to their own setting and data.

We believe this provides a framework that can guide how large language models can be used to adapt and deploy research pipelines in the future.

\section{Current Status of Project}

We have currently developed a research pipeline, focused on the specific question outlined above: the relationship between congressional debate about user data privacy and the innovation trajectory of technology companies, made available on the associated GitHub repository for this project. This pipeline demonstrates the transformation from raw transcripts of congressional hearings to a series of research outputs: textual summaries, summaries of perceived priorities (support/concerns identified by policymakers), and an index for policy relevance for publicly listed firms. This is a prototype of what can be done (using one set of approaches and design choices in the LLM pipeline). However, as we aim to show, and much like we are advocating, that this pipeline could be verified, adopted or adapted to other projects and applications.

\subsection{Expected Deliverables}

\begin{itemize}
    \item \textbf{Reproducible LLM Data Pipeline for Text Summarization, Stance Extraction, and Topic Relevance}
    \item \textbf{Public Datasets for Subsequent Research:} From the data of 36K Congressional Hearings from 1995--2025 (currently 2,352 hearings for the House and Senate Science Committees): full chunked text, intermediate summaries, list of topics discussed, summarized stances, an index of committee stances on topics discussed in hearings, and an index of relevance of those topics for publicly listed firms based on Item 1 business descriptions (1995--2025).
    \item \textbf{Illustrative Empirical Results:} Demonstration of the empirical results of the privacy regulation applied to investigate the impact of privacy regulation on the trajectory of corporate innovation.
    \item \textbf{Instructions and Framework for Deployment for Follow-on Research}
\end{itemize}

\subsection{Datasets (Outputs from Focal Research Project)}

\begin{itemize}
    \item \textbf{Text Summaries for Congressional Testimony from 1995 to 2025:} -- Intermediate outputs for verification and replication: Cleaned and chunked testimony, summarized testimony chunks, finalized text summaries.
    \item \textbf{Index of Committee Stances in Congressional Debate:} -- Intermediate outputs for verification and replication: Topics within congressional debate, filtered chunks focused with a stance on a focal topic, summarized positions on a focal topic, and an aggregated index of committee stances on topics discussed.
    \item \textbf{Index of Relevance of Congressional Debate for Publicly Listed Firms:} -- Intermediate outputs for verification and replication: Summarized congressional stances and company business descriptions based on 10-K filings, relevance index for individual company filings and congressional topic stances, and an aggregated index for topical relevance.
\end{itemize}

\subsection{Example of Empirical Application}

\begin{itemize}
    \item \textbf{Demonstration of the Policy Relevance Index:} Demonstrating the use of the constructed policy relevance index regarding the impact of congressional debate on the innovative trajectories of publicly listed firms. Using the context of the Cambridge Analytica scandal, we demonstrate how these tools can be used to identify firms that were in treated or control groups, allowing for an event study design to examine changes in innovation trajectories following these events.
\end{itemize}

\subsection{Key Risks and Limitations in Implementation}

We have tested the viability of our empirical strategy and have developed a working prototype of its deployment. We have confidence that this methodology will provide useful outputs. We have a goal of demonstrating this with as large a dataset as possible (the population of congressional testimony, public firm 10-Ks, and major policy topics discussed). We also recognize the potential for different ways of constructing and aggregating data. Because each of these approaches is time-consuming, we are committed to implementing all of them, but may become constrained in the number of individual specifications we can add.

\section{How This Advances Organizational and Social Science}

Long-form documents such as congressional hearings, regulatory filings, and corporate disclosures record how organizations and governments deliberate, yet they have been difficult to use: established text methods work best on short, information-dense text, and end-to-end LLM use relaxes this constraint only by introducing opacity. Our approach relaxes this constraint while alleviating the opacity of end-to-end use. If successful, it advances organizational research in three ways.

First, it provides a pipeline that turns long, noisy documents into measures whose construction can be inspected stage by stage. LLM-based measures thus become auditable research objects rather than outputs taken on trust, and the pipeline offers a template for reporting them, including intermediate outputs, stage-level validation, and stage-specific sensitivity analysis.

Second, it provides a public research dataset: summaries of more than 2,300 congressional hearings, a panel of legislators' stances by topic, and firm-level measures of exposure to policy debate. Because intermediate outputs are released alongside the final measures, other researchers can build on any stage, for example by re-aggregating stances by member or relating firm exposure to new outcomes.

Third, the framework enables a new form of replication. Because the pipeline is shared as a modular repository, researchers can direct coding agents to rebuild or adapt it at low cost. And because each stage produces outputs that can be compared with ours, a replication need not succeed or fail as a whole: discrepancies can be located at the stage where they arise. Replication thus shifts from an all-or-nothing comparison of final results to a stage-by-stage check.

Together, these contributions allow researchers both to do existing work better and to do new work. Existing work improves because the risks identified above become testable: when each stage is validated, systematic error can be detected and located before it biases a finding, and when a replication can localize discrepancies, disagreements over results can be traced to specific choices rather than left unresolved. New work becomes possible because researchers no longer face a choice between avoiding long-form documents and relying on measures that cannot be checked. Hearings, regulatory proceedings, and corporate disclosures, which record how organizations and governments deliberate, can support credible measures at scale; and because outputs are retained within documents, researchers can ask not only what a collective position is but who shapes it. The aim is not more research from the same materials, but better research from materials previously out of reach.

\vspace{2em}
\section*{References}
\begin{itemize}
    \item Asirvatham H, Mokski E, Shleifer A (2026) GPT as a measurement tool. NBER Working Paper, National Bureau of Economic Research, Cambridge, MA.
    \item Benoit K, De Marchi S, Laver C, Laver M, Ma J (2026) Using large language models to analyze political texts through natural language understanding. \textit{American Journal of Political Science} 70(4):1589--1605.
    \item Carlson NA, Burbano V (2026) The use of LLMs to annotate data in management research: Foundational guidelines and warnings. \textit{Strategic Management Journal} 47(3):699--725.
    \item Csaszar FA, Ketkar H, Kim H (2024) Artificial intelligence and strategic decision-making: Evidence from entrepreneurs and investors. \textit{Strategy Science} 9(4):322--345.
    \item de Kok T (2025) ChatGPT for textual analysis? How to use generative LLMs in accounting research. \textit{Management Science} 71(9):7888--7906.
    \item Dell'Acqua F, McFowland E, Mollick ER, Lifshitz-Assaf H, Kellogg K, Rajendran S, Krayer L, Candelon F, Lakhani KR (2026) Navigating the jagged technological frontier: Field experimental evidence of the effects of artificial intelligence on knowledge worker productivity and quality. \textit{Organization Science} 37(2):403--423.
    \item Dell M (2025) Deep learning for economists. \textit{Journal of Economic Literature} 63(1):5--58.
    \item Dell M, Rambachan A (2026) The measurement revolution? Credible measurement and inference in the age of AI. NBER Working Paper, National Bureau of Economic Research, Cambridge, MA.
    \item Doshi AR, Bell JJ, Mirzayev E, Vanneste BS (2025) Generative artificial intelligence and evaluating strategic decisions. \textit{Strategic Management Journal} 46(3):583--610.
    \item Gilardi F, Alizadeh M, Kubli M (2023) ChatGPT outperforms crowd workers for text-annotation tasks. \textit{Proceedings of the National Academy of Sciences} 120(30):e2305016120.
    \item Grimmer J, Stewart BM (2013) Text as data: The promise and pitfalls of automatic content analysis methods for political texts. \textit{Political Analysis} 21(3):267--297.
    \item Hannigan TR, Haans RF, Vakili K, Tchalian H, Glaser VL, Wang MS, Kaplan S, Jennings PD (2019) Topic modeling in management research: Rendering new theory from textual data. \textit{Academy of Management Annals} 13(2):586--632.
    \item Jha M, Qian J, Weber M, Yang B (2024) ChatGPT and corporate policies. NBER Working Paper, National Bureau of Economic Research, Cambridge, MA.
    \item Kim AG, Muhn M, Nikolaev VV (2023) Bloated disclosures: Can ChatGPT help investors process information? Chicago Booth Research Paper 23-07, University of Chicago, Booth School of Business, Chicago.
    \item Kogan L, Papanikolaou D, Seru A, Stoffman N (2017) Technological innovation, resource allocation, and growth. \textit{Quarterly Journal of Economics} 132(2):665--712.
    \item Rathje S, Mirea DM, Sucholutsky I, Marjieh R, Robertson CE, Van Bavel JJ (2024) GPT is an effective tool for multilingual psychological text analysis. \textit{Proceedings of the National Academy of Sciences} 121(34):e2308950121.
\end{itemize}


\clearpage
\appendix

\begin{center}
    {\huge \textsc{Technical Appendix}\par}
    \vspace{1.2em}
    {\Large \textsc{Opening the Black Box of Large Language Models in Text-based Measurement: Staged Pipelines with Verifiable Checkpoints}\par}
    \vspace{1em}
    {\normalsize\bfseries Milan Miric, Albert Roh, Nan Jia\par}
    \vspace{0.2em}
    {\small\itshape University of Southern California, Marshall School of Business\par}
    \vspace{0.8em}
    {\footnotesize\textbf{Proposal for Special Issue:} AI-Enabled Frontiers in Organizational Science (\textit{Organization Science})\par}
    \vspace{0.3em}
    {\footnotesize\textbf{Repository:} \url{https://github.com/HC-Geopolitics-Innovation/LLM-Text-Analysis-and-Synthesis}\par}
\end{center}

\vspace{1.5em}

\section*{Appendix A. Overview of Steps in Large Language Model Enabled Research Pipeline}
\addcontentsline{toc}{section}{Appendix A. Overview of Steps in Large Language Model Enabled Research Pipeline}

Here we summarize and outline the described pipeline, split across three related tasks (or problems) solved by LLMs.

\subsection*{A.1. Long-Document Text Summarization}

Many textual documents contain important information but also contain noisy and complex textual data. Large language models can provide a concise, information-rich summary that can be utilized in subsequent research tasks.

\subsubsection*{A.1.1. Overview of Processing Pipeline}

A pipeline to manage this process is as follows (summarized at length in Section xxx):

\begin{itemize}
    \item \textbf{Understand Document Structure and Define Semantic Partitioning:} Long textual documents often have a particular structure and therefore may be split into meaningful chunks. Context-aware chunking uses meaningful units of text in the document---such as a document section or a speaker's uninterrupted contribution---to split documents into chunks that can be processed by the LLM. The document is partitioned into these chunks and processed further.
    \item \textbf{Summarize and Combine the Inputs:} Each chunk of text is provided to an LLM with instructions outlining the desired format of the summary and the key information that should be extracted. The outputs can be combined in a variety of ways, such as recursively summarizing chunks or using a dedicated ``reduce'' stage to combine the summaries. This is referred to as a ``map-reduce'' approach.
    \item \textbf{Model Selection, Validation, and Deployment:} This is inherently an iterative process where researchers must test different approaches for partitioning, prompts for summarization, models for summarization, and combination strategies. All of these choices can be evaluated by inspecting intermediate outputs (partitions, summaries, outcomes for different models) and validated against ground truth where available. Once the optimal set of choices is identified, the pipeline is deployed across the entire corpus of documents.
\end{itemize}

\subsubsection*{A.1.2. Application: Summarizing Long Text Congressional Hearings}

We apply this pipeline to a corpus of U.S. Congressional hearing testimonies. The corpus consists of hearing transcripts spanning from 1995 to 2025, with multiple sessions per year. Hearings are typically long, often exceeding 100 pages, and feature multiple witnesses and committee members presenting testimony and engaging in Q\&A discussions.

To apply the pipeline, we first partition each hearing into individual witness testimonies. Within each testimony, we use sentence-level splitting to ensure granular coverage. These segments are then processed using the Map-Reduce summarization approach, with a ``reduce'' stage to consolidate summaries from multiple witnesses into a single comprehensive hearing summary.

These summaries are tested with different prompts, models, and splitting strategies. As a validation approach, we compare our generated summaries against human-generated summaries (abstracts) available on C-SPAN for these hearings. To identify the extent to which they are similar or comparable, we compute the cosine similarity between the embeddings of the human-written abstracts and the large language model generated summaries. As a way of providing a baseline for these values, in Figure 1, we compare them to the cosine similarity of different hearings from the same committee, or from hearings from different committees. Such tests allow us to confirm that the textual summaries are a useful distillation of the information in the long-form textual documents.

\subsubsection*{A.1.3. Research Outputs}

The summarization pipeline yields a transparent pipeline of code and prompts, as well as intermediate outputs at each stage of the process that can be utilized further:

\begin{itemize}
    \item \textbf{Partitioned Textual Documents:} A cleaned and partitioned version of each hearing that can be uploaded to an LLM without exceeding context windows.
    \item \textbf{Summarized Partitions:} A summary of each partition, generated by feeding the partitions into an LLM.
    \item \textbf{Committee Hearing Summaries:} Aggregated summaries from individual partitions into high-level summaries of overall hearing proceedings, capturing main themes, discussions, and conclusions (an output of the map-reduce process).
\end{itemize}

\subsection*{A.2. Extracting Topics, Positions, and Arguments from Long-Form Documents}

Large language models can also extract subtle but important nuances from textual documents, specifically related to topics of interest. In this step, we provide a method for extracting topics, positions, and arguments from long-form documents.

\subsubsection*{A.2.1. Overview of Processing Pipeline}

A pipeline to manage this process is as follows:

\begin{itemize}
    \item \textbf{Extract Topics from Textual Summaries:} From the long-form textual summaries, researchers can identify the major topics discussed within each summary. Because these are condensed, information-rich summaries, a variety of existing models can be effective for this task.
    \item \textbf{Filter Textual Summaries by Topic Relevance:} The partitioned text chunks contain subtle and nuanced arguments. However, as described in Section 3.1, many of these chunks contain noise, such as generic statements, remarks made by peripheral entities, or irrelevant information. Large language models can be used to filter out chunks that are relevant to having a stance on the focal topic.
    \item \textbf{Extract Arguments and Stances from Relevant Textual Summaries:} Within the filtered chunks, the LLM is tasked with extracting key positions related to the focal topic. This includes identifying whether the statement is supportive of or opposed to the topic, the specific nuances of the argument, and contextual information from the document.
    \item \textbf{Aggregate and Summarize:} The output of the preceding stage provides a set of concise data about the stances of individuals on a focal topic. These can be aggregated or summarized in different ways and at different levels of granularity.
    \item \textbf{Prototyping and Validation:} The goal of this exercise is to extract information that is otherwise difficult to capture and contextualize at scale. While this is challenging to validate quantitatively because comprehensive ground truth does not pre-exist, it is possible to compare against external baselines. We provide an example by validating our summaries against human-generated ones.
\end{itemize}

\subsubsection*{A.2.2. Application: Extracting Stances and Arguments around Focal Topics}

We apply this methodology to extract stances and arguments from U.S. Congressional hearings. First, we filter out summaries that are not relevant to the topic of interest and extract arguments and stances related to the focal topic. For instance, within congressional debate on a particular topic (such as the security and privacy of consumers), it becomes possible to capture both arguments related to the risks presented by a policy topic (e.g., potential misuse of data, impact on free speech) and potential opportunities (e.g., economic growth, innovation).

We aggregate these summaries in different ways, ultimately arriving at a numerical measure that captures the overall tenor of discussion within congressional debate. In doing so, we capture the predominant stances and arguments made by policymakers.

In Figure 2, we illustrate the output of this analysis across four interconnected panels (grouped in twos: Panels A \& B, Panels C \& D).

\subsubsection*{A.2.3. Research Outputs}

The outputs of this pipeline are:

\begin{itemize}
    \item \textbf{Topics of Debate within Congressional Hearings:} A standardized list of topics discussed within the textual documents, assigned to each hearing document.
    \item \textbf{Filtered Arguments and Examples on Each Topic:} A set of statements filtered from each congressional debate that captures the main arguments related to each specific topic, classified according to whether the policy is discussed as an opportunity or a threat.
    \item \textbf{Predominant Stances on Each Topic:} An index for each topic indicating the extent to which it is perceived as an opportunity or a threat by policymakers in each hearing, also aggregated at the year and committee levels to provide an overall description of main stances on each topic.
\end{itemize}

\subsection*{A.3. Policy Relevance for Publicly Listed Firms}

The earlier analysis provides detailed information extraction from congressional debate about economically important policies that directly shape the external environment of publicly listed firms. While there are anecdotal examples of firms being influenced by these policies, it is difficult to quantify this at scale and determine the extent to which firms are impacted.

To quantitatively measure the relevance of policy debate to publicly listed firms requires computing the relatedness between policy documents and corporate disclosures that reflect the business of the firm. This is fundamentally achieved by computing the textual distance between documents conveying policy relevance and documents describing the firm's business. LLMs can be used to either filter or summarize text, as described above, in order to improve the accuracy of this distance measure.

\subsubsection*{A.3.1. Overview of Processing Pipeline}

A pipeline to manage this process is as follows (summarized at length in Section xxx):

\begin{itemize}
    \item \textbf{Identify \& Pre-Process Focal Documents:} In the earlier process, we filtered out text chunks reflecting positions or stances of individuals or organizations on a particular topic. As an optional step, these can be further summarized into a stance summary on each issue. Similarly, an optional step is to summarize textual documents from companies (for example, Form 10-K filings).
    \item \textbf{Compute Semantic Distance and Topical Relevance:} Using either raw textual documents or summaries created in the earlier step, we compute the semantic relatedness of these documents to the 10-K statements of firms. This can be done using textual distance measures between document embeddings (e.g., cosine similarity using transformer-based embeddings) or by using large language models to compute semantic similarity. The robustness and sensitivity to different similarity measures can be tested and validated.
    \item \textbf{Aggregate Distance and Relevance Measures:} Outputs from the earlier phase can then be aggregated up to the level appropriate for a specific analysis. For instance, it may be relevant to compute the impact of a policy debate on a specific firm over time, a specific firm in aggregate, or across an entire industry. Fine-grained outputs at each stage of the process allow for validation, verification, and evaluation.
\end{itemize}

\subsubsection*{A.3.2. Application: Connecting Focal Concerns and Opportunities of Policymakers to Business Areas of Publicly Listed Firms}

We apply this methodology to understand how congressional policy debate around a particular policy topic impacts the business of firms. We focus on a particularly transformational shift in regulatory debate: the regulation of technology companies in relation to consumer privacy following the Cambridge Analytica scandal, which affected Facebook / Meta.

We utilize the filtered text from Section 3.2.3 to first create a summary of the Congressional debate on the topic of Privacy and Security, highlighting primary concerns and opportunities. We then compute the semantic distance between these summaries and the 10-K statements of publicly listed firms. We aggregate these measures to identify the firms whose businesses are potentially impacted by these policies. We report both raw data that can be leveraged for subsequent research and aggregate measures relevant to our empirical study.

We summarize the results of this classification in Figure 3. In Figure 3A (Panel A), we show the average cosine similarity between firms' Item 1 business descriptions and Congressional statements on Topic 68 (Consumer Data Privacy Legislation), relative to the sample mean across sectors. Firms in the Information (NAICS 51) sector exhibit the highest policy relevance, consistent with empirical expectations. In Figure 3B (Panel B), we show the top 10 firms most relevant to the Congressional debate on consumer data privacy legislation, finding that major digital and social media firms (such as Facebook/Meta and Alphabet/Google) emerge as the most exposed entities.

\subsubsection*{A.3.3. Research Outputs}

The outputs of this pipeline are:

\begin{itemize}
    \item \textbf{Summary of Policy Opportunities and Concerns for Each Congressional Session:} A textual summary of key concerns and opportunities related to a particular policy discussed by policymakers during a particular session.
    \item \textbf{Aggregated Summary of Opportunities and Concerns for Each Major Policy Issue:} A synthesized summary of primary opportunities and concerns related to each policy issue across policymakers.
    \item \textbf{Scores for Policy Relevance to Publicly Listed Firms:} A distance-based score measuring how much a particular policy (opportunity or threat) relates to the specific business activities of a publicly listed firm.
\end{itemize}

% =========================================================================
% FIGURES 1, 2, AND 3 AT THE END OF APPENDIX A
% =========================================================================

\clearpage
\begin{center}
\captionof{figure}{Cosine Similarity Between Hand-Coded Summaries and Large Language Model Generated Summaries}
\vspace{6pt}
\setlength{\fboxsep}{0pt}\setlength{\fboxrule}{0.4pt}%
\fbox{\includegraphics[width=\dimexpr\textwidth-2\fboxrule\relax]{"1 - Figure 1 - Cosine Similarity Between Hand-Coded and LLM Summaries.png"}}
\par\vspace{8pt}
{\figurenotesize\justifying\noindent\itshape Notes: Discriminant validity test comparing LLM-generated summaries against independent human-written C-SPAN abstracts using dense sentence embeddings (bge-large-en-v1.5). The target hearing exhibits high cosine similarity ($\rho=0.742$) against its own human abstract, substantially exceeding baselines from different hearings within the same committee ($\rho=0.518$) and hearings from different committees ($\rho=0.384$).\par}
\end{center}

\clearpage
\begin{center}
\captionof{figure}{Topics and Overview of Stance (Support / Concern of Policymakers)}
\vspace{4pt}
{\small\textbf{Panel A: Distribution of Member Net Stances Across Congressional Hearings}}\par\vspace{2pt}
\setlength{\fboxsep}{0pt}\setlength{\fboxrule}{0.4pt}%
\fbox{\includegraphics[width=\dimexpr\textwidth-2\fboxrule\relax,height=0.38\textheight,keepaspectratio]{"2 - Figure 2 Panel A - Distribution of Member Net Stances.png"}}\par\vspace{6pt}
{\small\textbf{Panel B: Longitudinal Net Stance by Hearing Over Time (1998--2024)}}\par\vspace{2pt}
\fbox{\includegraphics[width=\dimexpr\textwidth-2\fboxrule\relax,height=0.38\textheight,keepaspectratio]{"3 - Figure 2 Panel B - Net Stance by Hearing Over Time.png"}}
\end{center}

\clearpage
\begin{center}
\addtocounter{figure}{-1}
\captionof{figure}{Topics and Overview of Stance (Continued)}
\vspace{4pt}
{\small\textbf{Panel C: Subtopic Polarization Pre- vs. Post-2016 (Cambridge Analytica)}}\par\vspace{2pt}
\setlength{\fboxsep}{0pt}\setlength{\fboxrule}{0.4pt}%
\fbox{\includegraphics[width=\dimexpr\textwidth-2\fboxrule\relax,height=0.36\textheight,keepaspectratio]{"4 - Figure 2 Panel C - Subtopic Polarization Pre-Post 2016.png"}}\par\vspace{6pt}
{\small\textbf{Panel D: Polarizing vs. Consensual Subtopics Across Parties}}\par\vspace{2pt}
\fbox{\includegraphics[width=\dimexpr\textwidth-2\fboxrule\relax,height=0.36\textheight,keepaspectratio]{"5 - Figure 2 Panel D - Polarizing vs Consensual Subtopics.png"}}
\par\vspace{8pt}
{\figurenotesize\justifying\noindent\itshape Notes: In Panel A, we report the distribution of topics discussed across the hearing sample and the balance of stances in support of or concern regarding these topics. In Panel B, we report the breakdown for topics specifically related to privacy and security. From Panel A, it is evident that consumer privacy is one of the dominant topics, and the subgroup Consumer Data Privacy Legislation is among the few where substantial legislative support emerges for regulatory action. In Panels C and D, we report the longitudinal share of stances discussed by Congress over time, separating Consumer Data Privacy Legislation into its own group (Panel D) and contrasting it with other technology policy topics (Panel C). We find a pronounced structural increase in the share of statements supporting consumer privacy legislation after 2016 (following the Cambridge Analytica disclosures). No corresponding increase is observed among other policy topics. This analysis demonstrates how the pipeline documents political debate and traces legislative shifts over time.\par}
\end{center}

\clearpage
\begin{center}
\captionof{figure}{Topic Relevance Scores for Publicly Listed Firms}
\vspace{3pt}
{\small\textbf{Panel A: NAICS Sector Cosine Relative to Sample Mean}}\par\vspace{2pt}
\setlength{\fboxsep}{0pt}\setlength{\fboxrule}{0.4pt}%
\fbox{\includegraphics[width=\dimexpr\textwidth-2\fboxrule\relax,height=0.36\textheight,keepaspectratio]{"2 - Figure 3 Panel A - Sector Relevance (NAICS 2-Digit).png"}}\par\vspace{4pt}
{\small\textbf{Panel B: Top 10 Relevant Firms by Share of Sentences}}\par\vspace{2pt}
\fbox{\includegraphics[width=\dimexpr\textwidth-2\fboxrule\relax,height=0.36\textheight,keepaspectratio]{"3 - Figure 3 Panel B - Top 10 Relevant Firms.png"}}
\par\vspace{8pt}
{\figurenotesize\justifying\noindent\itshape Notes: Panel A reports firm summary cosine by NAICS sector relative to the sample mean (Topic 68 policy summary vector). Each bar represents the sector mean minus the mean across all 841 firms in the three-year sample (fiscal years 2010, 2015, and 2020) of the firm relevance score, computed as the cosine similarity between the firm’s Item 1 summary and the policy summary of Topic 68 hearing statements. Zero denotes the sample mean (0.5143); bars to the right indicate sectors above the mean, and bars to the left indicate sectors below. Error bars span one standard deviation of firm scores within each sector, capturing within-industry cross-firm variation. Sectors reflect two-digit NAICS codes, with manufacturing (31--33), retail trade (44--45), and transportation/warehousing (48--49) grouped. Sectors with fewer than 5 firms are excluded. Panel B reports the top 10 most relevant firms by share of sentences aligned with the Topic 68 statement centroid. Bars report the share of each firm’s Item 1 sentences whose cosine similarity with the Topic 68 policy statement centroid is at or above 0.75, divided by total Item 1 sentences. Ranked among 841 scored firms across fiscal years 2010, 2015, and 2020. Least relevant firms are omitted as 679 firms tie at zero exposure.\par}
\end{center}

% =========================================================================
% APPENDIX B
% =========================================================================

\clearpage
\section*{Appendix B. Overview of Datasets Used for Data Pipeline}
\addcontentsline{toc}{section}{Appendix B. Overview of Datasets Used for Data Pipeline}

This appendix provides an integrated overview of the primary input corpora, verification benchmarks, and research output datasets utilized across our AI-enabled pipeline.

\subsection*{A. Primary Data Inputs}

\textbf{Congressional Hearings Dataset (U.S. GPO GovInfo CHRG Collection):}
\begin{itemize}
    \item \textbf{Enumerated Universe:} Complete package manifest of 36,378 hearings spanning 1995--2025, enumerated prior to data retrieval to ensure every omitted item is a measured subtraction rather than an unobserved omission.
    \item \textbf{Focal Science Sample:} For this proposal, we focus on hearings before the House Science, Space, and Technology Committee (hssy00) and the Senate Commerce, Science, and Transportation Committee (sscm00), comprising 2,352 hearings containing 93.6 million words across 116,000 speaker turns.
    \item \textbf{Data Processing \& Integrity:} Downloaded MODS XML metadata and full HTML transcripts parsed into speaker-attributed speaking turns, witness lists, member rosters, and hearing dates. Systematic defect tracking identified 1,253 un-digitized stubs (predominantly from 1995--1996) and 4,448 empty member rosters, both flagged in indicator columns rather than silently dropped.
\end{itemize}

\textbf{Publicly Listed Firms Business Descriptions (SEC Form 10-K Filings):}
\begin{itemize}
    \item \textbf{Enumerated Archive:} Full SEC EDGAR Form 10-K archive accessed via sec-api.io, comprising 305,140 annual reports across 394 monthly containers (34 GB of text spanning 1993--2025).
    \item \textbf{Firm Linking \& Identification:} Mapped to financial fundamentals through a bijective CRSP--Compustat crosswalk linking 25,207 CIK $\leftrightarrow$ GVKEY identifier pairs.
    \item \textbf{Item 1 Extraction \& Verification:} The Business description section (Item 1) is extracted for all S\&P 500 reports. The Extractor API returned 432 filings (1.9\%) empty; all 432 were deterministically recovered and verified from the local bulk archive.
    \item \textbf{Usable Corporate Sample:} Yields 22,299 S\&P 500 annual reports complete with extracted text, provenance metadata, and usability flags, directly feeding Stage 4 policy relevance scoring.
\end{itemize}

\subsection*{B. Verification Data}

\textbf{Independent Reference Abstracts (C-SPAN Archives API):}
\begin{itemize}
    \item \textbf{Scope:} 18,697 committee programs covering 1995--2025, retrieved from api.c-spanarchives.org.
    \item \textbf{Content \& Role:} Human-authored editorial abstracts (median 51 words) written by professional C-SPAN editors independently of congressional staff and our pipeline.
    \item \textbf{Pipeline Function:} Serves as an unpolluted, external ground-truth benchmark for discriminant validity testing, comparing cosine similarity of LLM summaries against human abstracts across identical versus distinct committees.
\end{itemize}

\subsection*{C. Output Research Datasets}

\begin{itemize}
    \item \textbf{Text Summarization Corpus (Stage 2 Outputs):} Cleaned, turn-attributed, and chunked transcript segments across all 2,352 hearings; intermediate chunk-level summaries produced by turn-aware map-reduce prompts; finalized, entity-attributed committee hearing summaries (\texttt{hearing\_summaries.parquet}; 2,352 records) validated against independent C-SPAN abstracts.
    \item \textbf{Policy Stance Panel Dataset (Stage 3 Outputs):} Ratified, canonical 111-topic policy taxonomy codebook (\texttt{frozen\_codebook\_111.parquet}) extracted from congressional testimony and crosswalked to the Comparative Agendas Project (CAP); filtered text statements isolating discrete member policy positions (support, concern, neutral); longitudinal congressional member-topic stance panel (\texttt{member\_stance\_panel.parquet}) tracking policy orientation across committees and congresses over time.
    \item \textbf{Corporate Topic Relevance Dataset (Stage 4 Outputs):} Synthesized Form 10-K Item 1 corporate business description embeddings across 22,299 firm-year filings; synthesized congressional debate summaries and topic centroid vectors (e.g., Topic 68: Consumer Data Privacy Legislation); firm-year policy exposure and relevance score panel (\texttt{firm\_topic\_relevance\_panel.parquet}) linking S\&P 500 firms to congressional debate dynamics.
\end{itemize}

% =========================================================================
% APPENDIX C
% =========================================================================

\clearpage
\section*{Appendix C. Example of Research Enabled by LLM Data Pipeline}
\addcontentsline{toc}{section}{Appendix C. Example of Research Enabled by LLM Data Pipeline}

\subsection*{Evaluating the Impact of Privacy Policy Debate Following the Cambridge Analytica Scandal on the Trajectory of Firm Corporate Research on Privacy Technology}

As described above, we demonstrate how the application of our data pipeline to congressional testimony and company 10-K filings allows us to quantify which firms were impacted by this policy and study this question in a way not previously possible.

Leveraging the rich data pipeline we describe above, we are able to generate an index which provides us with an indication for the stance of a policy debate (e.g., related to corporate / customer privacy; see Section A.3.2), but also simultaneously the extent to which these policy positions would have impacted firms, and which firms would have been impacted. Crucially, these issues would not have been possible to identify, let alone empirically study, in the absence of an approach such as the one we propose, because while it would have been possible to identify prominent technology companies from the data, it would not have been easy to conclude which are the companies which were directly impacted by these policy changes, beyond the usual suspects (e.g., Meta / Facebook following the Cambridge Analytica scandal).

In Figure 4, Panel A, we plot the distribution of our generated topic relevance scores, indicating each quartile, and then for each corresponding quartile compute the intensity of patenting by firms in that quartile in the domain of user privacy and data security. The firms that we identify as more likely to be influenced by the privacy policy debate are also those generating an overwhelming share of patents in privacy-related domains. We report the breakdown across a variety of specific subclasses related to privacy and technologies involved in collecting and processing customer data.

In Panel B, we report the change in the annual rate of privacy patents filed by firms within each topic relevance quartile before and after 2016 (2009--2015 average vs. 2016--2020 average), coinciding with the structural surge in Congressional debate regarding technology and consumer privacy legislation (see Section A.2.2). As shown in Panel B, the increase in corporate privacy innovation is overwhelmingly concentrated among firms in Quartile 4 (the highest topic relevance firms), where annual patent filings rose from an average of 5.57 to 8.84 patents per firm per year (+3.27 patents/year/firm), while lower-relevance firms exhibited negligible increases.

In Panel C, we investigate the specific privacy technology subclasses directly and compare how annual patent filings shifted across the pre- and post-2016 periods. These subclasses capture distinct technological approaches to collecting, tracking, and managing customer data. The patterns in Panel C indicate that the post-2016 surge was particularly driven by an increase in companies anonymizing or de-correlating personal data from user identifiers (\texttt{G06F 21/6254}, \texttt{H04W 12/033}), which rose from 27.9 to 121.2 patents per year (+335.1\%). Simultaneously, there was a dramatic emergence of alternative network-level behavioral tracking technologies (\texttt{H04L 67/535}), which jumped from 0.7 to 37.0 patents per year (+5,080\%), allowing firms to infer user attributes from online activity rather than direct profile data. Importantly, this reflects a shift toward different technological approaches from those that were predominant prior to 2016.

In Figure 5, we decompose these results further by investigating how these patterns appear across patent filings over time. In Panel A, we plot the annual count of distinct patents filed by sample firms in each topic group across the sample period (Tracking User Profiles and Data, Protecting Privacy and Anonymity in PII, Anonymizing or Decorrelating Data, and Protecting Confidentiality by Encryption). We observe a large increase in patents related to tracking user profiles and collecting customer data, particularly increasing after 2009, coinciding with the widespread adoption of smartphones. This is mirrored by an increase (albeit less pronounced) of technology designed to protect customer data. Not reflected in this figure, but in this period, with the transmission of digital data through digital mobile networks, saw an increase in patenting of patents related to encryption and data protection over mobile networks by hardware and infrastructure manufacturers.

We see a gradual increase in the patenting of anonymization and data de-correlation technologies accelerating after 2016, mirrored with a decline in the tracking of user data. This coincides with the increase in government scrutiny and broader concerns around privacy and legislation during that period. Following 2018, we similarly observe a decline in the protection of customer data and data de-correlation technologies, along with data collection in general. This coincides with the Cambridge Analytica scandal as well as the introduction of GDPR. Following this policy however we observe an increase in a number of other technologies, namely the tracking of users through network and interface activities (how and what they search) as well as through technologies that protect confidentiality through encryption.

In Panel B, we show the patenting within these categories by various major technology companies. We can observe that these are the major technology companies patenting within this space, but also can note meaningful variation between technology companies. For instance, hardware manufacturers such as Apple, Qualcomm, Cisco and Intel are more likely to patent in technologies related to encryption, whereas software and digital service providers invested heavily in customer data tracking and profiling technologies. Moreover, it is notable that while Meta and Adobe invested heavily in collecting customer data, they invested considerably less in protecting confidentiality or anonymizing and de-correlating data than peer service companies (e.g. Alphabet or Amazon).

In Figure 6, we investigate this further by exploring how the probability of firms pursuing these technologies changed in relation to 2015. This is by no means an exogenous event. The increase in policy discussion coincided with a broader debate and pressure on how companies would have the rights to manage customer data. However, it provides a useful instance in which we can evaluate the usefulness of this index.

In Panel A, we report the coefficients for the probability of firms patenting in the domain of Anonymizing or De-correlating technologies, where we observed a sizeable increase following 2016 in the raw descriptive data. We observe that those firms in the top quartile, had been patenting and continued to patent in this space (Q4). By contrast, firms in the lower quartiles (Q1--Q2) did not nor did they adjust their patenting behaviour. However, firms in Q3, which were somewhat at risk of increasing their patenting behaviour experienced an increase in patenting following 2015, which had been previously flat. However, following 2018, those same firms decreased their patenting, while the upper quartiles continued to patent these technologies. This illustrates the heterogeneity of how the innovation output of these firms changed alongside this change in the regulatory debate. Moreover, following 2018 the divergence between the first quartile and the other quartiles increased further. In Panel B, we focus on those technologies which were beginning to emerge and then became accelerated following the introduction of GDPR. In this case we observe how the upper quartile (Q4) of companies increased their patenting behaviour in these technologies as compared to the other companies including those in Q3 who had clearly engaged in other privacy technologies, but had not made this pivot.

In summary, this empirical exercise demonstrates how the LLM enabled text analysis pipeline can create a framework which can be used to generate measures which would otherwise be infeasible. Using these measures we are able to generate results and insights which can enable research such as this into the innovation response of firms following changes in the policy environment and legislative debate.

% =========================================================================
% FIGURES 4, 5, AND 6 AT THE END OF APPENDIX C (EACH ON ITS OWN PAGE)
% =========================================================================

\clearpage
\begin{center}
\captionof{figure}{Firm Topic Relevance and Privacy Patenting Pre- and Post-2016}
\vspace{4pt}
{\small\textbf{Panel A: Distribution of Firm Topic Relevance and Privacy Patents Across Categories and Quartiles}}\par\vspace{2pt}
\setlength{\fboxsep}{0pt}\setlength{\fboxrule}{0.4pt}%
\fbox{\includegraphics[width=\dimexpr\textwidth-2\fboxrule\relax,height=0.45\textheight,keepaspectratio]{"1 - Figure 4 Panel A - Distribution of Firm Topic Relevance and Privacy Patents.png"}}
\end{center}

\clearpage
\begin{center}
\addtocounter{figure}{-1}
\captionof{figure}{Firm Topic Relevance and Privacy Patenting Pre- and Post-2016 (Continued)}
\vspace{3pt}
{\small\textbf{Panel B: Annual Privacy Patent Filings Pre- vs. Post-2016 by Relevance Quartile}}\par\vspace{2pt}
\setlength{\fboxsep}{0pt}\setlength{\fboxrule}{0.4pt}%
\fbox{\includegraphics[width=\dimexpr\textwidth-2\fboxrule\relax,height=0.36\textheight,keepaspectratio]{"2 - Figure 4 Panel B - Annual Privacy Patent Filings by Relevance Quartile.png"}}\par\vspace{4pt}
{\small\textbf{Panel C: Annual Privacy Patent Filings Pre- vs. Post-2016 by Subclass}}\par\vspace{2pt}
\fbox{\includegraphics[width=\dimexpr\textwidth-2\fboxrule\relax,height=0.36\textheight,keepaspectratio]{"3 - Figure 4 Panel C - Annual Privacy Patent Filings by Subclass.png"}}
\par\vspace{8pt}
{\figurenotesize\justifying\noindent\itshape Notes: Figure 4 presents firm-level topic relevance and corporate innovation responses across 1,058 publicly traded sample firms (excluding Facebook/Meta and Alphabet/Google). Panel A plots the distribution of firm topic relevance (cosine similarity between firm Form 10-K Item 1 business descriptions and the Congressional policy summary for Topic 68: Consumer Data Privacy Legislation; left axis, scaled 0--100, dodger blue bars without outlines). Red dotted vertical lines denote quartile boundaries ($Q_1 \le 0.486$, $Q_2$: 0.486--0.511, $Q_3$: 0.511--0.537, $Q_4 > 0.537$). Right Axis 1 (solid black line, scale 0--125) traces average total privacy patents per firm across 1999--2021 (rising from 2.2 in $Q_1$ to 111.1 in $Q_4$). Right Axis 2 (scale 0--12, black spine) traces average patents per firm across the five narrow technology subclasses: Tracking Profiles \& Targeted Ads (G06Q 30/0269, 30/0255, 30/0261; navy solid line, $Q_4$: 8.6), Protecting Privacy \& PII (H04W 12/02; dark grey dashed line, $Q_4$: 5.2), Anonymizing Data (G06F 21/6254, H04W 12/033; medium grey dash-dot line, $Q_4$: 3.0), Tracking User Activity (H04L 67/535; navy dotted line, $Q_4$: 1.2), and Protecting Confidentiality Encryption (H04W 12/03; light grey loosely dashed line, $Q_4$: 0.5). Panel B reports annual privacy patents filed per firm before (2009--2015 annual mean, black) and after (2016--2020 annual mean, dodger blue) by topic relevance quartile, with capless standard error bars. Panel C reports annual patent filings across the five technology subclasses before and after 2016 (black = Pre-2016, dodger blue = Post-2016).\par}
\end{center}

\clearpage
\begin{center}
\captionof{figure}{Long-Term Trends and Firm Concentration in Privacy Patenting (1999--2021)}
\vspace{4pt}
{\small\textbf{Panel A: Annual Privacy Patents by Topic Group (1999--2021)}}\par\vspace{2pt}
\setlength{\fboxsep}{0pt}\setlength{\fboxrule}{0.4pt}%
\fbox{\includegraphics[width=\dimexpr\textwidth-2\fboxrule\relax,height=0.36\textheight,keepaspectratio]{"1 - Figure 5 Panel A - Aggregate Filing Trends across Privacy CPC Codes.png"}}\par\vspace{4pt}
{\small\textbf{Panel B: Top 20 Privacy Patenting Firms Split by Technology Group}}\par\vspace{2pt}
\fbox{\includegraphics[width=\dimexpr\textwidth-2\fboxrule\relax,height=0.36\textheight,keepaspectratio]{"2 - Figure 5 Panel B - Top 20 Patent Assignees by Privacy Technology Group.png"}}
\par\vspace{8pt}
{\figurenotesize\justifying\noindent\itshape Notes: Figure 5 illustrates the longitudinal evolution and corporate concentration of privacy patents across the sample of 1,058 publicly traded firms from 1999 to 2021. The sample is constructed by matching Compustat firm identifiers (GVKEY) to USPTO utility patent grants using cooperative patent classification (CPC) codes. Panel A plots annual distinct patent filings across the four focal topic groups: Tracking User Profiles and Data ($N = 2,764$), Protecting Privacy and Anonymity in PII ($N = 1,519$), Anonymizing or Decorrelating Data ($N = 954$), and Protecting Confidentiality by Encryption ($N = 142$). Faint vertical dashed lines indicate regulatory milestones: 2009.5 (FTC privacy reports and HITECH Act), 2016.5 (Cambridge Analytica disclosures and GDPR adoption), and 2018.5 (GDPR enforcement and CCPA enactment). Panel B displays the 20 sample firms with the highest total privacy patent holdings over 1999--2021, with stacked bars indicating the compositional share of each firm's portfolio across the focal technology classes.\par}
\end{center}

\clearpage
\begin{center}
\captionof{figure}{Event Study of Patenting Probability by Topic Relevance Quartile}
\vspace{4pt}
{\small\textbf{Panel A: Anonymizing or Decorrelating Technologies}}\par\vspace{2pt}
\setlength{\fboxsep}{0pt}\setlength{\fboxrule}{0.4pt}%
\fbox{\includegraphics[width=\dimexpr\textwidth-2\fboxrule\relax,height=0.34\textheight,keepaspectratio]{"1 - Figure 6 Panel A - Event Study of Patenting Probability - Anonymizing Data.png"}}\par\vspace{4pt}
{\small\textbf{Panel B: Tracking User Activity Technologies}}\par\vspace{2pt}
\fbox{\includegraphics[width=\dimexpr\textwidth-2\fboxrule\relax,height=0.34\textheight,keepaspectratio]{"2 - Figure 6 Panel B - Event Study of Patenting Probability - Tracking User Activity.png"}}
\par\vspace{8pt}
{\figurenotesize\justifying\noindent\itshape Notes: Figure 6 presents event-study estimates tracking the annual probability that a firm files at least one patent within designated privacy technology subclasses from 2009 to 2020 relative to 2015 across topic relevance quartiles ($Q_1$ to $Q_4$). The sample comprises 1,058 publicly traded firms. The outcome in each panel is an indicator variable equal to one if firm $i$ filed at least one patent application in year $t$ in the focal technology group, estimated using firm fixed effects with standard errors clustered at the firm level. Vertical capless error bars represent 95\% confidence intervals. Vertical red dashed background lines mark the 2015/2016 policy transition (pre/post-2016 split) and the 2018/2019 GDPR enforcement and CCPA enactment milestone. The horizontal dashed line denotes zero change relative to 2015. Panel A evaluates Anonymizing or Decorrelating Data (CPC codes G06F 21/6254, H04W 12/033). Pre-2016 coefficients are flat and indistinguishable from zero across all quartiles. Following the 2015/2016 transition, firms in the highest relevance quartile ($Q_4$, dodger blue) exhibit a pronounced surge in patenting probability, increasing by +6.0 percentage points in 2019 ($p < 0.001$) and +6.4 percentage points in 2020 ($p < 0.001$), while lower quartiles ($Q_1$ through $Q_3$) remain flat. Panel B evaluates Tracking User Activity (CPC code H04L 67/535). Pre-trends from 2009 to 2015 remain essentially zero across all quartiles. Following the 2015/2016 transition, $Q_4$ firms experience a continuous, monotonic rise in tracking patenting probability: +1.5 percentage points in 2016, +2.3 percentage points in 2017, +3.4 percentage points in 2018, +5.7 percentage points in 2019, and +9.8 percentage points in 2020 ($p < 0.001$). Quartiles $Q_1$ and $Q_2$ exhibit no significant change, while $Q_3$ registers a modest uptick to +1.9 percentage points in 2020.\par}
\end{center}

% =========================================================================
% APPENDIX D
% =========================================================================

\clearpage
\section*{Appendix D. Reproducible and Adaptable Data Pipelines with Coding Agents}
\addcontentsline{toc}{section}{Appendix D. Reproducible and Adaptable Data Pipelines with Coding Agents}

While this data pipeline leverages AI tools to perform important analysis tasks, an additional important benefit is emerging in empirical research: research papers made available through platforms such as GitHub in easy machine-readable formats (Markdown, Python code, Stata, and R) can be deployed rapidly through coding agents to build entire data pipelines. In addition to providing a methodology that can advance research with text data, we also advocate for and provide a framework for using online repositories of papers as a medium where researchers can share tools that are easily deployable and adaptable for other scholars. To demonstrate this in the accompanying repository folder (\href{https://github.com/HC-Geopolitics-Innovation/LLM-Text-Analysis-and-Synthesis/tree/main/X2%20-%20Example%20of%20Reproducing%20and%20Adapting%20the%20Research%20Pipeline}{\texttt{X2 - Example of Reproducing and Adapting the Research Pipeline/}}), we provide a sample agent prompt (\href{https://github.com/HC-Geopolitics-Innovation/LLM-Text-Analysis-and-Synthesis/blob/main/X2%20-%20Example%20of%20Reproducing%20and%20Adapting%20the%20Research%20Pipeline/X1%20-%20Sample%20Agent%20Prompt%20For%20Adapting%20Research%20Pipeline%20To%20New%20Question.md}{\texttt{X1 - Sample Agent Prompt For Adapting Research Pipeline To New Question.md}}) and an end-to-end implementation plan generated by an LLM (\href{https://github.com/HC-Geopolitics-Innovation/LLM-Text-Analysis-and-Synthesis/blob/main/X2%20-%20Example%20of%20Reproducing%20and%20Adapting%20the%20Research%20Pipeline/X2%20-%20Sample%20Implementation%20Plan%20Generated%20by%20Pipeline.md}{\texttt{X2 - Sample Implementation Plan Generated by Pipeline.md}}) demonstrating how to adapt this methodology to evaluate corporate innovation responses to Congressional debates over renewable energy and carbon emissions. Researchers can do the same with their own AI tools, augment the implementation plan, and deploy it further.

\end{document}
